\documentclass[12pt,a4paper]{ctexbook}
\usepackage{geometry}
\geometry{margin=2.2cm}
\usepackage{setspace}
\setstretch{1.35}
\usepackage{hyperref}
\title{花间集·huajianji-1-juan}
\author{古典文献整理}
\date{}
\begin{document}
\maketitle
\tableofcontents
\newpage
\section{菩萨蛮 其一}
\textbf{温庭筠}\par\vspace{0.5em}
小山重叠金明灭，鬓云欲度香腮雪。\par
懒起画蛾眉，弄妆梳洗迟。\par
照花前后镜，花面交相映。\par
新帖绣罗襦，双双金鹧鸪。\par

\section{菩萨蛮 其二}
\textbf{温庭筠}\par\vspace{0.5em}
水晶帘里玻璃枕，暖香惹梦鸳鸯锦。\par
江上柳如烟，雁飞残月天。\par
藕丝秋色浅，人胜参差剪。\par
双鬓隔香红，玉钗头上风。\par

\section{菩萨蛮 其三}
\textbf{温庭筠}\par\vspace{0.5em}
蕊黄无限当山额，宿妆隐笑纱窗隔。\par
相见牡丹时，暂来还别离。\par
翠钗金作股，钗上蝶双舞。\par
心事竟谁知，月明花满枝。\par

\section{菩萨蛮 其四}
\textbf{温庭筠}\par\vspace{0.5em}
翠翘金缕双鸂鶒，水纹细起春池碧。\par
池上海棠梨，雨晴红满枝。\par
绣衫遮笑靥，烟草粘飞蝶。\par
青琐对芳菲，玉关音信稀。\par

\section{菩萨蛮 其五}
\textbf{温庭筠}\par\vspace{0.5em}
杏花含露团香雪，绿杨陌上多离别。\par
灯在月胧明，觉来闻晓莺。\par
玉钩褰翠幕，妆浅旧眉薄。\par
春梦正关情，镜中蝉鬓轻。\par

\section{菩萨蛮 其六}
\textbf{温庭筠}\par\vspace{0.5em}
玉楼明月长相忆，柳丝袅娜春无力。\par
门外草萋萋，送君闻马嘶。\par
画罗金翡翠，香烛消成泪。\par
花落子规啼，绿窗残梦迷。\par

\section{菩萨蛮 其七}
\textbf{温庭筠}\par\vspace{0.5em}
凤凰相对盘金缕，牡丹一夜经微雨。\par
明镜照新妆，鬓轻双脸长。\par
画楼相望久，栏外垂丝柳。\par
音信不归来，社前双燕回。\par

\section{菩萨蛮 其八}
\textbf{温庭筠}\par\vspace{0.5em}
牡丹花谢莺声歇，绿杨满院中庭月。\par
相忆梦难成，背窗灯半明。\par
翠钿金压脸，寂寞香闺掩。\par
人远泪阑干，燕飞春又残。\par

\section{菩萨蛮 其九}
\textbf{温庭筠}\par\vspace{0.5em}
满宫明月梨花白，故人万里关山隔。\par
金雁一双飞，泪痕沾绣衣。\par
小园芳草绿，家住越溪曲。\par
杨柳色依依，燕归君不归。\par

\section{菩萨蛮 其十}
\textbf{温庭筠}\par\vspace{0.5em}
宝函钿雀金鸂鶒，沉香阁上吴山碧。\par
杨柳又如丝，驿桥春雨时。\par
画楼音信断，芳草江南岸。\par
鸾镜与花枝，此情谁得知。\par

\section{菩萨蛮 其十一}
\textbf{温庭筠}\par\vspace{0.5em}
南园满地堆轻絮，愁闻一霎清明雨。\par
雨后却斜陽，杏花零落香。\par
无言匀睡脸，枕上屏山掩。\par
时节欲黄昏，无聊独倚门。\par

\section{菩萨蛮 其十二}
\textbf{温庭筠}\par\vspace{0.5em}
夜来皓月才当午，重帘悄悄无人语。\par
深处麝烟长，卧时留薄妆。\par
当年还自惜，往事那堪忆。\par
花落月明残，锦衾知晓寒。\par

\section{菩萨蛮 其十三}
\textbf{温庭筠}\par\vspace{0.5em}
雨暗夜合玲珑日，万枝香袅红丝拂。\par
闲梦忆金堂，满庭萱草长。\par
绣帘垂簏簌，眉黛远山绿。\par
春水渡溪桥，凭栏魂欲消。\par

\section{菩萨蛮 其十四}
\textbf{温庭筠}\par\vspace{0.5em}
竹风轻动庭除冷，珠帘月上玲珑影。\par
山枕隐秾妆，绿檀金凤凰。\par
两蛾愁黛浅，故国吴宫远。\par
春恨正关情，画楼残点声。\par

\section{更漏子 其一}
\textbf{温庭筠}\par\vspace{0.5em}
柳丝长，春雨细，花外漏声迢递。\par
惊塞雁，起城乌，画屏金鹧鸪。\par
香雾薄，透帘幕，惆怅谢家池阁。\par
红烛背，绣帘垂，梦长君不知。\par

\section{更漏子 其二}
\textbf{温庭筠}\par\vspace{0.5em}
星斗稀，钟鼓歇，帘外晓莺残月。\par
兰露重，柳风斜，满庭堆落花。\par
虚阁上，倚栏望，还似去年惆怅。\par
春欲暮，思无穷，旧欢如梦中。\par

\section{更漏子 其三}
\textbf{温庭筠}\par\vspace{0.5em}
金雀钗，红粉面，花里暂时相见。\par
如我意，感君怜，此情须问天。\par
香作穗，蜡成泪，还似两人心意。\par
山枕腻，锦衾寒，觉来更漏残。\par

\section{更漏子 其四}
\textbf{温庭筠}\par\vspace{0.5em}
相见稀，相忆久，眉浅淡烟如柳。\par
垂翠幕，结同心，待郎熏绣衾。\par
城上月，白如雪，蝉鬓美人愁绝。\par
宫树暗，鹊桥横，玉签初报明。\par

\section{更漏子 其五}
\textbf{温庭筠}\par\vspace{0.5em}
背江楼，临海月，城上角声呜咽。\par
堤柳动，岛烟昏，两行征雁分。\par
京口路，归帆渡，正是芳菲欲度。\par
银烛尽，玉绳低，一声村落鸡。\par

\section{更漏子 其六}
\textbf{温庭筠}\par\vspace{0.5em}
玉炉香，红蜡泪，偏照画堂秋思。\par
眉黛薄，鬓云残，夜长衾枕寒。\par
梧桐树，三更雨，不道离情正苦。\par
一叶叶，一声声，空阶滴到明。\par

\section{归国遥·香玉}
\textbf{温庭筠}\par\vspace{0.5em}
香玉，翠凤宝钗垂簏簌，胜金粟，越罗春水绿。\par
画堂照帘残烛，梦余更漏促。\par
谢娘无限心曲，晓屏山断续。\par

\section{归国遥·双脸}
\textbf{温庭筠}\par\vspace{0.5em}
双脸，小凤战蓖金飐艳。舞衣无力风敛，藕丝秋色染。\par
锦帐绣帷斜掩，露珠清晓簟。\par
粉心黄蕊花靥，黛眉山两点。\par

\section{酒泉子·花映柳条}
\textbf{温庭筠}\par\vspace{0.5em}
花映柳条，闲向绿萍池上。\par
凭栏干，窥细浪，雨萧萧。\par
近来音信两疏索，洞房空寂寞。\par
掩银屏，垂翠箔，度春宵。\par

\section{酒泉子·日映纱窗}
\textbf{温庭筠}\par\vspace{0.5em}
日映纱窗，金鸭小屏山碧。\par
故乡春，烟霭隔，背兰釭。\par
宿妆惆怅倚高阁，千里云影薄。\par
草初齐，花又落，燕双双。\par

\section{酒泉子·楚女不归}
\textbf{温庭筠}\par\vspace{0.5em}
楚女不归，楼枕小河春水。\par
月孤明，风又起，杏花稀。\par
玉钗斜簪云鬟髻，裙上金缕凤。\par
八行书，千里梦，雁南归。\par

\section{酒泉子·罗带惹香}
\textbf{温庭筠}\par\vspace{0.5em}
罗带惹香，犹系别时红豆。\par
泪痕新，金缕旧，断离肠。\par
一双娇燕语雕梁，还是去年时节。\par
绿陰浓，芳草歇，柳花狂。\par

\section{定西番·汉使昔年离别}
\textbf{温庭筠}\par\vspace{0.5em}
汉使昔年离别。\par
攀弱柳，折寒梅，上高台。\par
千里玉关春雪，雁来人不来。\par
羌笛一声愁绝，月徘徊。\par

\section{定西番·海燕欲飞调羽}
\textbf{温庭筠}\par\vspace{0.5em}
海燕欲飞调羽。萱草绿，杏花红，隔帘栊。\par
双鬓翠霞金缕，一枝春艳浓。\par
楼上月明三五，琐窗中。\par

\section{定西番·细雨晓莺春晚}
\textbf{温庭筠}\par\vspace{0.5em}
细雨晓莺春晚。\par
人似玉，柳如眉，正相思。\par
罗幕翠帘初卷，镜中花一枝。\par
肠断塞门消息，雁来稀。\par

\section{杨柳枝·宜春苑外最长条}
\textbf{温庭筠}\par\vspace{0.5em}
宜春苑外最长条，闲袅春风伴舞腰。\par
正是玉人肠绝处，一渠春水赤栏桥。\par

\section{杨柳枝·南内墙东御路傍}
\textbf{温庭筠}\par\vspace{0.5em}
南内墙东御路傍，须知春色柳丝黄。\par
杏花未肯无情思，何事行人最断肠。\par

\section{杨柳枝·苏小门前柳万条}
\textbf{温庭筠}\par\vspace{0.5em}
苏小门前柳万条，毵毵金线拂平桥。\par
黄莺不语东风起，深闭朱门伴舞腰。\par

\section{杨柳枝·金缕毵毵碧瓦沟}
\textbf{温庭筠}\par\vspace{0.5em}
金缕毵毵碧瓦沟，六宫眉黛惹香愁。\par
晚来更带龙池雨，半拂栏干半入楼。\par

\section{杨柳枝·馆娃宫外邺城西}
\textbf{温庭筠}\par\vspace{0.5em}
馆娃宫外邺城西，远映征帆近拂提。\par
系得王孙归意切，不关芳草绿萋萋。\par

\section{杨柳枝·两两黄鹂色似金}
\textbf{温庭筠}\par\vspace{0.5em}
两两黄鹂色似金，袅枝啼露动芳音。\par
春来幸自长如线，可惜牵缠荡子心。\par

\section{杨柳枝·御柳如丝映九重}
\textbf{温庭筠}\par\vspace{0.5em}
御柳如丝映九重，凤凰窗映绣芙蓉。\par
景陽楼畔千条路，一面新妆等晓风。\par

\section{杨柳枝·织锦机边莺语频}
\textbf{温庭筠}\par\vspace{0.5em}
织锦机边莺语频，停梭垂泪忆征人。\par
塞门三月犹萧索，纵有垂杨未觉春。\par

\section{南歌子·手里金鹦鹉}
\textbf{温庭筠}\par\vspace{0.5em}
手里金鹦鹉，胸前绣凤凰。\par
偷眼暗形相，不如从嫁与，作鸳鸯。\par

\section{南歌子·似带如丝柳}
\textbf{温庭筠}\par\vspace{0.5em}
似带如丝柳，团酥握雪花。\par
帘卷玉钩斜，九衢尘欲暮，逐香车。\par

\section{南歌子·鬓堕低梳髻}
\textbf{温庭筠}\par\vspace{0.5em}
鬓堕低梳髻，连娟细扫眉。\par
终日两相思，为君憔悴尽，百花时。\par

\section{南歌子·脸上金霞细}
\textbf{温庭筠}\par\vspace{0.5em}
脸上金霞细，眉间翠钿深。\par
倚枕覆鸳衾，隔帘莺百啭，感君心。\par

\section{南歌子·扑蕊添黄子}
\textbf{温庭筠}\par\vspace{0.5em}
扑蕊添黄子，呵花满翠鬟。\par
鸳枕映屏山，明月三五夜，对芳颜。\par

\section{南歌子·转盼如波眼}
\textbf{温庭筠}\par\vspace{0.5em}
转盼如波眼，娉婷似柳腰。\par
花里暗相招，忆君肠欲断，恨春宵。\par

\section{南歌子·懒拂鸳鸯枕}
\textbf{温庭筠}\par\vspace{0.5em}
懒拂鸳鸯枕，休缝翡翠裙。\par
罗帐罢炉熏，近来心更切，为思君。\par

\section{河渎神·河上望丛柯}
\textbf{温庭筠}\par\vspace{0.5em}
河上望丛柯，庙前春雨来时。\par
楚山无限鸟飞迟，兰棹空伤别离。\par
何处杜鹃啼不歇，艳红开尽如血。\par
蝉鬓美人愁绝，百花芳草佳节。\par

\section{河渎神·孤庙对寒潮}
\textbf{温庭筠}\par\vspace{0.5em}
孤庙对寒潮，西陵风雨萧萧。\par
谢娘惆怅倚兰桡，泪流玉箸千条。\par
暮天愁听思归乐，早梅香满山郭。\par
回首两情萧索，离魂何处飘泊。\par

\section{河渎神·铜鼓赛神来}
\textbf{温庭筠}\par\vspace{0.5em}
铜鼓赛神来，满庭幡盖徘徊。\par
水村江浦过风雷，楚山如画烟开。\par
离别橹声空萧索，玉容惆怅妆薄。\par
青麦燕飞落落，卷帘愁对珠阁。\par

\section{女冠子·含娇含笑}
\textbf{温庭筠}\par\vspace{0.5em}
含娇含笑，宿翠残红窈窕。\par
鬓如蝉，寒玉簪秋水，轻纱卷碧烟。\par
雪胸鸾镜里，琪树凤楼前。\par
寄语青娥伴，早求仙。\par

\section{女冠子·霞帔云发}
\textbf{温庭筠}\par\vspace{0.5em}
霞帔云发，钿镜仙容似雪。\par
画愁眉，遮语回轻扇，含羞下绣帷。\par
玉楼相望久，花洞恨来迟。\par
早晚乘鸾去，莫相遗。\par

\section{玉蝴蝶·霞帔云发}
\textbf{温庭筠}\par\vspace{0.5em}
秋风凄切伤离，行客未归时。\par
塞外草先衰江，南雁到迟。\par
芙蓉凋嫩脸，杨柳堕新眉。\par
摇落使人悲，断肠谁得知。\par

\end{document}